\begin{description}
   \item[dataブロック] データ長L，各時点の体重Wがデータです。加えて，推定すべき欠損値の数もデータとして受け取ります。
   \item[parametersブロック] データ長と同じだけの状態を表す\verb|mu|と，状態の変動幅\verb|tau|，状態に付加される計測誤差の変動幅\verb|sig|を宣言しています。また最初の状態$\mu_0$を表すパラメタと，欠損の数だけある推定用パラメータ\verb|Miss_W|を用意しました。
   \item[modelブロック] 先ほど同様，0時点目の状態から出てくるものですので，$\mu_1 \sim N(\mu_0,\tau)$をモデル化しています。次に各回のデータについて，$W_l$が999でなければ(同じでない，という論理式は\verb|!=|で表現します)，普通の状態空間モデルのように尤度として計算します。 そうではない，すなわちデータ$W_l=999$である，つまり欠損値であるはずということなら，$W_{miss} \sim N(\mu_l,\sigma)$で推定してやります。ここで特殊な変数$j$が出てきています。これは「何番目の欠損値か」を表す変数です。ブロックの中の一部だけで変数を宣言するため，まず\verb|for|文全体を中括弧で括り，はじめに$j=0$という数字を宣言しています。ここで欠損値推定のシーンになると，$j$のカウターを1つ増加させ，$j$番目の欠損値の推定をさせる，というやり方をしています。欠損値の数と行番号が合致しないので，こうした工夫が必要なのですね。その後のコードは先ほどと同じです。
\end{description}
